\section{Related Work and Literature Survey}
\label{sec:related}
Computational drug discovery and structural biology have undergone rapid transformation between 2021 and 2026, transitioning from rigid grid-based docking protocols to continuous deep generative diffusion, all-atom foundation co-folding, and sparse neural routing architectures.

\subsection{Biomolecular Co-Folding and Structure Foundation Models}
The paradigm of macromolecular structure determination was reshaped by AlphaFold 2 and subsequent multi-chain assemblies. Google DeepMind's AlphaFold 3 \cite{abramson2024accurate} unified multi-entity modeling by replacing invariant representations with a pairformer diffusion trunk that predicts joint coordinates of proteins, nucleic acids, ions, and small molecules directly from sequences and SMILES. Concurrently, Wohlwend et al.\ introduced Boltz-1 \cite{wohlwend2024boltz1}, establishing open-source all-atom biomolecular modeling by matching AlphaFold 3 accuracy across protein--protein and protein--nucleic acid benchmarks while substantially reducing inference latency. In late 2026, Geodesic Intelligence announced NovaAtom-Lite \cite{geodesic2026novaatom}, setting a state-of-the-art benchmark on the FoldBench evaluation suite (scoring 79.07\% on 172 antibody--antigen complexes) through hardware-optimized continuous geometric representations. Parallel architectures, including IsoDDE \cite{fedde2023isodde}, SeedFold \cite{fedoseev2024seedfold}, and Protenix-v1 \cite{zheng2024protenix}, have established structural fidelity across complex multimeric assemblies.

Despite their structural accuracy, these co-folding networks are inherently evaluative rather than de novo generative. They predict the 3D equilibrium conformation of pre-existing, user-supplied chemical structures. They cannot natively sample, generate, or optimize novel, drug-like small-molecule scaffolds from empty target binding pockets, nor do they model the evolutionary trajectory of acquired mutational drug resistance.

\subsection{Equivariant Generative SBDD and Continuous Flow Matching}
Structure-based drug design (SBDD) has developed through equivariant graph neural networks and score-based 3D diffusion. DiffSBDD \cite{schneuing2022diffsbdd} and TargetDiff \cite{guan2023targetdiff} pioneered continuous 3D atom coordinate diffusion conditioned on protein pocket residues, enforcing rotational and translational equivariance under the $\mathrm{SE}(3)$ Lie group via equivariant graph convolutional networks (EGNNs). In parallel, Pocket2Mol \cite{peng2022pocket2mol} demonstrated efficient auto-regressive pocket sampling. More recently, continuous-time Riemannian flow matching \cite{lipman2023flow,albergo2023building} has emerged as an alternative to discrete diffusion, providing straight-line probability trajectories and reducing sampling steps.

However, standard generative diffusion models suffer from severe physical and chemical limitations:
\begin{enumerate}
    \item \textbf{High Steric Clash and Bond Distortion Rates}: Buttenschoen et al.\ demonstrated via the PoseBusters benchmark \cite{van2023posebusters} that over 45\% of poses generated by TargetDiff and Pocket2Mol violate physical stereochemical constraints, exhibiting overlapping atomic van der Waals radii, distorted aromatic planarity, and non-physical tetrahedral bond angles.
    \item \textbf{Single-Target Overfitting}: Generative models typically sample ligands within a single rigid crystallographic cavity, yielding candidates that fail when the target acquires a single-point mutation in clinical settings.
    \item \textbf{Decoupled Biophysical and Pharmacokinetic Objectives}: Existing models generate molecules without evaluating systemic clearance, human diploid cell safety margins, or quantum frontier orbital kinetics during generation.
\end{enumerate}

\subsection{Sparse Mixture-of-Experts (MoE) in Scientific Deep Learning}
Sparse Mixture-of-Experts architectures, originally formalized by Shazeer et al.\ \cite{shazeer2017outrageously} and expanded in large-scale language models such as Mixtral \cite{jiang2024mixtral}, decouple total parameter capacity from per-step computational cost by routing input tokens to a sparse subset of specialized sub-networks. In computational biophysics, molecular design requires simultaneous optimization across disparate biophysical regimes: quantum electronic charge transfers, macromolecular shape complementarity, steric clash penalties, and pharmacokinetic clearance. Monolithic dense networks suffer from negative gradient interference across these conflicting objectives. By allocating dedicated experts to affinity, geometry, safety, resistance, quantum kinetics, and retrosynthesis, sparse MoE enables uncompromised multi-objective Pareto optimization.

\subsection{Acquired Gatekeeper Resistance and Structural Escape Corridors}
Targeted oncology therapeutics face inevitable resistance driven by Darwinian selective pressure. The primary mechanism involves secondary gatekeeper mutations that alter the stereoelectronic landscape of the binding pocket:
\begin{itemize}
    \item \textbf{KRAS G12C and Switch-II Pocket Escape}: The first-in-class covalent KRAS G12C inhibitor Sotorasib binds within the cryptic switch-II groove \cite{ostrem2013kras,canon2019clinical}. Acquired secondary substitutions, notably Y96D and A146T, introduce severe steric repulsion and electrostatic perturbation, abolishing drug binding while preserving oncogenic signaling \cite{tanaka2021clinical}.
    \item \textbf{EGFR C797S Gatekeeper Mutation}: 3rd-generation covalent EGFR inhibitors (e.g., Osimertinib) target the nucleophilic Cys797 thiol. The emergent tertiary C797S substitution replaces cysteine with serine, abolishing covalent bond formation and reducing potency by three orders of magnitude \cite{thress2015acquired}.
    \item \textbf{BCR-ABL1 T315I Gatekeeper Mutation}: In chronic myeloid leukemia, the T315I substitution replaces threonine with bulky isoleucine, abolishing critical hydrogen bonding with frontline kinase inhibitors \cite{shah2004overriding}. While Ponatinib overrides T315I, its potent inhibition of VEGFR2/KDR induces fatal arterial thrombotic events in up to 30\% of treated patients.
    \item \textbf{ESR1 Y537S Conformational Shift}: In metastatic hormone receptor-positive breast cancer, the Y537S substitution stabilizes the active transcription factor conformation of the estrogen receptor independent of estradiol, rendering traditional selective estrogen receptor modulators ineffective.
    \item \textbf{PfDHFR Quadruple Escape}: In infectious disease, Plasmodium falciparum dihydrofolate reductase acquires a quadruple mutation (N51I/C59R/S108N/I164L) that confers multi-thousand-fold resistance to Pyrimethamine \cite{singh2011resurgence}.
\end{itemize}

\subsection{Algorithmic Retrosynthesis and Green Chemistry}
A common pitfall of generative de novo drug discovery is the generation of synthetically intractable molecular structures. Modern computer-aided synthesis planning tools, such as AiZynthFinder \cite{genheden2020aizynthfinder}, evaluate disconnections against cataloged reaction templates and commercial building block libraries (e.g., Enamine, Sigma-Aldrich). Furthermore, international regulatory standards (ICH Q2/Q14) and green chemistry frameworks mandate low Environmental Factors ($E\text{-Factor} < 10$), high atom economy, and scalable convergent routes for viable preclinical development.

AETHER-RAMI bridges these disparate domains by integrating continuous-time Riemannian flow matching on $\mathrm{SE}(3)$ with an 8-expert sparse MoE network, enforcing stereochemical PoseBusters Hamiltonian constraints, and embedding clinical resistance mutational landscapes directly into the generative objective.